% R2-Dreamer DMC-Vision walker_walk -> walker_run transfer — experiment report.
%
% Figures (reports/figures/r2dreamer_*.png) are the exact committed outputs of
% notebook 06 (06_r2dreamer_results.ipynb), extracted from its rendered cells.
%
% Build: tectonic r2dreamer_report.tex   (or pdflatex, twice)

\documentclass[11pt]{article}
\usepackage[margin=2.6cm]{geometry}
\usepackage{graphicx}
\usepackage{booktabs}
\usepackage{amsmath}
\usepackage{xcolor}
\usepackage[colorlinks=true, linkcolor=blue!50!black, urlcolor=blue!50!black, citecolor=blue!50!black]{hyperref}
\graphicspath{{figures/}}
\emergencystretch=3em  % absorb residual long-typewriter-token line breaks

\title{Representation-Regularized World Models for Visual Locomotion
       Transfer:\\ An R2-Dreamer \texttt{walker\_walk}\,$\rightarrow$\,\texttt{walker\_run}
       Study on DMC Vision}
\author{Thomas Georges\thanks{Engineering and experiment automation developed
        with Claude (Anthropic) in the \texttt{wm-prediction} repository.}}
\date{June 16, 2026}

\begin{document}
\maketitle

\begin{abstract}
We study transfer in a reconstruction-free, representation-regularized
DreamerV3 variant (``R2-Dreamer''~\cite{r2dreamer,dreamerv3}) on the DeepMind
Control Suite~\cite{dmc} from pixels. A 25M-parameter world model is trained on
\texttt{walker\_walk} and then either fine-tuned on \texttt{walker\_run} from the
walk checkpoint, or trained on \texttt{walker\_run} from scratch, under an
identical 400k-step target budget. The source policy effectively solves the walk
task (eval return $968$). On the run task, fine-tuning reaches $742$ versus $563$
from scratch, a $+179$ ($+32\%$) advantage at equal budget, with a much larger
head start earlier: the warm-started policy scores $\sim\!375$ before any
run-specific training (vs.\ $\sim\!25$ for scratch) and leads by roughly
$4\times$ at the 100k-step mark, with the gap peaking near $+448$ and then
narrowing as the scratch baseline---still rising at the budget limit---catches
up. Training is stable for all three runs, and the R2 redundancy-reduction
objective converges cleanly. A PCA analysis of the recurrent latent states shows
that the two same-task policies occupy clearly separated regions of feature
space and that the fine-tuned model traces a larger, cyclically structured
(limit-cycle-like) latent manifold, while the scratch model collapses into a
compact blob---so fine-tuning changes not only learning speed but the learned
representation itself. Results are single-seed and at a fixed budget; we discuss
the resulting limits on the conclusions.
\end{abstract}

\section{Objective}

The goal of this track is a faithful, budget-feasible demonstration of the
canonical Dreamer transfer experiment---train a latent world model on a source
task, save and reload it, adapt it to a shifted task, and compare fine-tuning
against training from scratch---using the \emph{representation-regularized}
R2-Dreamer objective on image observations. Concretely we run the three-way
comparison
\[
\begin{aligned}
\texttt{source\_base} &:~\texttt{walker\_walk}~\text{(from scratch)},\\
\texttt{target\_finetune} &:~\texttt{walker\_run}~\text{(initialized from \texttt{source\_base})},\\
\texttt{target\_scratch} &:~\texttt{walker\_run}~\text{(from scratch)},
\end{aligned}
\]
all on \texttt{dmc\_vision} ($64\times64$ RGB), orchestrated from a single Colab
notebook against an unmodified upstream checkout,\footnote{\texttt{NM512/r2dreamer}~\cite{r2dreamer},
itself a PyTorch reimplementation of DreamerV3~\cite{dreamerv3} extending
\texttt{NM512/dreamerv3-torch}~\cite{dreamerv3torch}; behavioural changes are
applied as marker-guarded, idempotent, reversible textual patches at launch
time (checkpoint load/save, interval checkpoints, headless-render guard,
serial-env fallback).} with all checkpoints, metrics, videos, and figures
persisted to Google Drive. Every number and figure in this report is read
directly from the committed outputs of the results notebook
(\texttt{notebooks/06\_r2dreamer\_results.ipynb}).

\section{Model}
\label{sec:model}

\paragraph{Backbone.} The agent is a DreamerV3-style latent world
model~\cite{dreamerv3,dreamerv2,dreamer}. A convolutional encoder maps each
$64\times64\times3$ frame to an embedding; a Recurrent State-Space Model
(RSSM)~\cite{planet} carries a deterministic recurrent state and a stochastic
latent and predicts forward dynamics; reward and episode-continuation heads
decode task quantities from the latent; and an actor--critic is trained
entirely ``in imagination'' on rollouts of the learned dynamics rather than on
environment frames~\cite{dreamer}. The observation here is vision only (the
encoder receives \texttt{\{image:(64,64,3)\}} and no proprioceptive MLP input).

\paragraph{Representation objective (the ``R2'' in R2-Dreamer).} The upstream
code exposes a swappable representation loss
(\texttt{model.rep\_loss} $\in$ \{\texttt{dreamer}, \texttt{r2dreamer},
\texttt{infonce}, \texttt{dreamerpro}\}). The standard \texttt{dreamer} setting
learns the latent through
pixel reconstruction with a decoder. We use \texttt{rep\_loss=r2dreamer}: a
\emph{reconstruction-free} objective that learns the representation through a
projector head trained with a redundancy-reduction (Barlow-Twins-style) loss
on projected latents~\cite{barlow}, placing it in the same reconstruction-free
model-based family as DreamerPro~\cite{dreamerpro} and contrastive
CPC/InfoNCE~\cite{cpc} alternatives. A practical consequence is that the
\texttt{r2dreamer} model has \emph{no decoder}: there is no imagined-frame
``\texttt{video\_pred}'' path (which the upstream agent exposes only for
\texttt{rep\_loss=dreamer}). Visual analysis therefore uses environment-rendered
policy rollouts and the geometry of the RSSM latent states, not decoded
reconstructions.

\paragraph{Parameter budget.} The nominal \texttt{size25M} configuration
instantiates $26.53$M trainable parameters, dominated by the RSSM and the R2
projector head:

\begin{table}[h]
\centering
\small
\begin{tabular}{lrl}
\toprule
Component & Parameters & Role \\
\midrule
RSSM                 & 13{,}442{,}304 & latent dynamics (recurrent + stochastic state) \\
Projector            &  5{,}898{,}240 & R2 representation head (redundancy reduction) \\
Value (critic)       &  1{,}869{,}951 & imagined-return value \\
Actor (policy)       &  1{,}776{,}396 & imagination policy \\
Reward predictor     &  1{,}573{,}503 & reward head \\
Continue predictor   &  1{,}475{,}713 & episode-continuation head \\
Encoder (CNN)        &    493{,}824 & $64\times64\times3 \rightarrow$ embedding \\
\midrule
\textbf{Total}       & \textbf{26{,}529{,}931} & \\
\bottomrule
\end{tabular}
\caption{World-model components and parameter counts, read from the rollout/
extraction logs. Each saved checkpoint is $325.9$\,MB (fp32 weights plus Adam
optimizer moments); the six most recent interval checkpoints are retained per
run (\texttt{checkpoint\_keep=6}).}
\label{tab:model}
\end{table}

\section{Experimental setup and resources}
\label{sec:setup}

\paragraph{Tasks and protocol.} Source task \texttt{dmc\_walker\_walk}, target
task \texttt{dmc\_walker\_run}, both rendered as $64\times64$ RGB
(Figure~\ref{fig:obs}). The fine-tune run initializes \emph{all} world-model
weights from the source checkpoint with a strict state-dict load and does
\emph{not} carry over optimizer state (\texttt{pretrained\_strict=true},
\texttt{load\_optimizer=false}); the scratch run is identical except for
random initialization. Both target runs share the same 400k-step budget so the
comparison is at equal target experience.

\begin{table}[h]
\centering
\small
\begin{tabular}{ll@{\hspace{1.6em}}ll}
\toprule
Setting & Value & Setting & Value \\
\midrule
Observation        & \texttt{dmc\_vision} $64{\times}64{\times}3$ & Train ratio        & 224 \\
Model              & \texttt{size25M}                            & Parallel envs      & 8 \\
Representation loss & \texttt{r2dreamer} (Barlow-style)          & Eval cadence       & every 50k steps, 5 episodes \\
Batch              & 16 sequences $\times$ length 64             & Seed               & 0 \\
Source steps       & 800k (reached $\sim$796k)                   & Target steps       & 400k (reached $\sim$396k) \\
\bottomrule
\end{tabular}
\caption{Training configuration (the \texttt{a100\_r2\_vision25m} three-way
preset in \texttt{configs/r2dreamer/}).}
\label{tab:config}
\end{table}

\paragraph{Compute and software.} Training and post-hoc analysis ran on a single
NVIDIA~A100 (Google Colab). The stack is PyTorch with Hydra configuration
(\texttt{NM512/r2dreamer}~\cite{r2dreamer} over
\texttt{dreamerv3-torch}~\cite{dreamerv3torch}), DeepMind Control / \texttt{dm\_control}
tasks~\cite{dmc,dmcontrol} on the MuJoCo physics engine~\cite{mujoco}, on a
Python~3.11 runtime (the version upstream declares). Image observations are
produced by headless CPU rendering through OSMesa
(\texttt{MUJOCO\_GL=osmesa}) while the policy runs on CUDA; a render
sanity-check confirms non-blank frames before any rollout
(Figure~\ref{fig:obs}). The full run writes model-only interval checkpoints at
evaluation boundaries; the source run in this report was interrupted and
resumed (three appended metric segments, 326 rows discarded), and the analysis
uses the final continuous segment to $\sim$796k steps.

\begin{figure}[h]
\centering

\includegraphics[width=0.16\linewidth]{r2dreamer_render_check.png}
\caption{A representative $64\times64$ RGB observation from the patched DMC
vision environment (render sanity-check: non-blank, pixel mean $68$, std $33$).
The policies learn locomotion from frames like this; because
\texttt{rep\_loss=r2dreamer} is decoder-free, the model never reconstructs
them.}
\label{fig:obs}
\end{figure}

\section{Results}
\label{sec:results}

\begin{table}[h]
\centering
\small
\begin{tabular}{llrrrr}
\toprule
Run & Init & Task & Final eval & Final train & Env steps \\
\midrule
\texttt{source\_base}     & random            & \texttt{walker\_walk} & $968.3$ & $964.2$ & $\sim$796k \\
\texttt{target\_finetune} & \texttt{source\_base} & \texttt{walker\_run}  & $742.3$ & $693.4$ & $\sim$396k \\
\texttt{target\_scratch}  & random            & \texttt{walker\_run}  & $563.1$ & $529.7$ & $\sim$396k \\
\bottomrule
\end{tabular}
\caption{Final results (DMC episodic return, $0$--$1000$). ``Eval'' is the
deterministic 5-episode evaluation score; ``train'' is the rolling training
return. For all three runs the best eval score equals the final eval score,
i.e.\ no late-training collapse.}
\label{tab:results}
\end{table}

\subsection{Source task: \texttt{walker\_walk} is effectively solved}

The source model reaches an evaluation return of $968.3$ on
\texttt{walker\_walk} (Table~\ref{tab:results}; the DMC walk reward saturates
near $1000$). This is the prerequisite for a meaningful transfer test: the
checkpoint we fine-tune from is a competent walker, not a half-trained one.
Figure~\ref{fig:source} shows the monotone climb to that ceiling.

\begin{figure}[h]
\centering
\includegraphics[width=0.62\linewidth]{r2dreamer_source_eval_curve.png}
\caption{Source-task evaluation return over $\sim$796k environment steps. The
\texttt{walker\_walk} model converges to $\approx\!968$, effectively solving the
task and providing the initialization for the fine-tune run.}
\label{fig:source}
\end{figure}

\subsection{Transfer: fine-tuning dominates scratch at every budget}

This is the central result (Figure~\ref{fig:transfer}). Three features stand
out. \emph{(i) A large warm start:} at the first evaluation, before any
run-specific gradient steps, the fine-tuned policy already returns
$\sim\!375$ on \texttt{walker\_run}, versus $\sim\!25$ for the random scratch
policy---walk and run share enough locomotion structure that the source model
transfers almost zero-shot. \emph{(ii) Faster convergence:} by 100k target
steps the fine-tune is at $\sim\!600$ while scratch is near $\sim\!150$, a
$\sim\!4\times$ gap; the fine-tune passes $700$ by $\sim$300k. \emph{(iii) A
persistent but shrinking lead:} the fine-tune ends at $742$ and scratch at
$563$, a $+179$ ($+32\%$) advantage at the shared 400k budget.

\begin{figure}[h]
\centering
\begin{minipage}{0.49\linewidth}
  \includegraphics[width=\linewidth]{r2dreamer_target_eval_curve.png}
\end{minipage}\hfill
\begin{minipage}{0.49\linewidth}
  \includegraphics[width=\linewidth]{r2dreamer_finetune_advantage.png}
\end{minipage}
\caption{Target-task transfer. \textbf{Left:} evaluation return on
\texttt{walker\_run} for fine-tune (green) vs.\ scratch (red); fine-tune leads
at every evaluation point. \textbf{Right:} the step-aligned advantage
(fine-tune $-$ scratch), which starts at $\sim\!352$, peaks at $\sim\!448$
around 100k steps, and then declines steadily to $+179$ at the budget limit as
the scratch baseline---still rising and not yet plateaued---closes part of the
gap.}
\label{fig:transfer}
\end{figure}

The advantage curve makes the dynamics explicit: transfer buys the most
\emph{early} (peak gap near 100--150k steps), and the steady decline thereafter
indicates that the result demonstrates \emph{sample efficiency}, not
necessarily a higher asymptotic ceiling---the scratch curve had not flattened
at 400k (Section~\ref{sec:caveats}).

\subsection{Training dynamics and optimization health}

The training-episode curves (Figure~\ref{fig:train}) corroborate the evaluation
story and reveal a brief adaptation transient: the fine-tune begins around
$\sim\!400$ (the inherited walk competence), \emph{dips} to $\sim\!270$ as the
walk gait is repurposed for running---a short, recoverable negative-transfer
phase---then climbs steadily past the scratch run for the remainder of
training. The scratch run rises roughly linearly from $\sim\!30$ and is still
ascending at the budget limit.

\begin{figure}[h]
\centering
\includegraphics[width=0.72\linewidth]{r2dreamer_training_episode.png}
\caption{Rolling-mean training episode return (window 8). Source
\texttt{walker\_walk} (blue) climbs to $\sim\!950$ over 800k steps; the
\texttt{walker\_run} fine-tune (green) starts high, dips during gait
re-adaptation, and converges near $\sim\!690$; scratch (red) rises from near
zero to $\sim\!530$ and is still increasing.}
\label{fig:train}
\end{figure}

Optimization is stable for all three runs (Figure~\ref{fig:diag}): the total
optimizer loss and the dynamics/KL, representation, value, and policy terms all
descend and remain bounded, with no divergence. In particular the
R2-specific Barlow redundancy-reduction term drops sharply and flattens,
indicating the reconstruction-free representation objective trained as intended.

\begin{figure}[h]
\centering
\includegraphics[width=0.86\linewidth]{r2dreamer_diagnostics.png}
\caption{Optimization diagnostics across the three runs (source blue, fine-tune
green, scratch red). All loss components are bounded and convergent; the
\texttt{train/loss/barlow} panel (the R2 redundancy-reduction term) converges
cleanly. The source run is noisier at high step counts but does not diverge.}
\label{fig:diag}
\end{figure}

\subsection{Latent geometry: the two policies live in different basins}

We extracted RSSM recurrent features along three evaluation episodes
($1500$ states each) for the fine-tune and scratch \texttt{walker\_run}
policies and projected them with PCA (Figure~\ref{fig:pca}). Although both
policies solve the \emph{same} task, their latent states occupy clearly
separated regions---near-linearly separable along the first principal
component---so the warm-started model encodes the task differently from the
scratch model rather than merely reaching the same representation faster. The 3D
view is more striking: the fine-tuned model's latents trace a large, organized
ring/limit-cycle manifold consistent with a periodic running gait, whereas the
scratch model's latents collapse into a small, compact cluster, mirroring its
lower return and less dynamic behaviour.

\begin{figure}[h]
\centering
\begin{minipage}{0.42\linewidth}
  \includegraphics[width=\linewidth]{r2dreamer_latent_pca_2d.png}
\end{minipage}\hfill
\begin{minipage}{0.55\linewidth}
  \includegraphics[width=\linewidth]{r2dreamer_latent_pca_3d.png}
\end{minipage}
\caption{Shared-basis PCA of RSSM latent trajectories, fine-tune (blue) vs.\
scratch (red) on \texttt{walker\_run}. \textbf{Left (2D):} the two policies
separate along PC1. \textbf{Right (3D):} the fine-tune traces a large,
cyclically structured manifold; the scratch policy occupies a compact blob.
The top three principal components explain only $\approx\!9.5/3.8/3.5\%$ of the
variance ($\sim\!17\%$ total), so these projections capture the dominant linear
structure of a high-dimensional latent space, not its entirety
(Section~\ref{sec:caveats}).}
\label{fig:pca}
\end{figure}

\subsection{Rollout videos}

The results notebook renders one camera rollout per checkpoint
(\texttt{source\_base}, \texttt{target\_finetune}, \texttt{target\_scratch})
and a rotating 3D latent-PCA animation, all embedded inline. Because the
\texttt{r2dreamer} model is decoder-free, these are true environment renderings
of the executed policy, not imagined reconstructions. Their qualitative content
tracks the scores in Table~\ref{tab:results}: a near-solved walk for the source,
a competent run for the fine-tune, and a weaker run for scratch. (A minor
corroborating detail: the smoothest, most periodic motion---the source
walk---produces the smallest encoded MP4.)

\section{Discussion}

\paragraph{What the experiment shows.} On a visual locomotion shift, a
representation-regularized DreamerV3 world model transfers strongly: a competent
\texttt{walker\_walk} model gives an almost zero-shot head start on
\texttt{walker\_run} ($\sim\!375$ vs.\ $\sim\!25$), learns several times faster
mid-training ($\sim\!4\times$ at 100k steps), and retains a clear $+179$
advantage at a fixed 400k budget. The whole machinery---strict checkpoint
save/reload across tasks, the reconstruction-free R2 objective, and headless
vision rendering---works end to end, and the comparison is clean because source,
fine-tune, and scratch differ only in initialization.

\paragraph{Efficiency versus ceiling.} The shrinking advantage curve
(Figure~\ref{fig:transfer}, right) is the honest qualifier: scratch was still
improving at the budget limit, so the evidence supports a large
\emph{sample-efficiency} benefit from transfer rather than a proven higher
ceiling. Whether the $+179$ gap is permanent or merely a head start that a much
longer scratch run would erase is not resolved by this budget.

\paragraph{Representation, not just speed.} The latent-geometry result is the
most interesting qualitative finding. Two policies trained to the same task from
different initializations end up in measurably different, near-linearly
separable representational regions, and the better policy carries a richer,
cyclically structured latent manifold. This suggests fine-tuning lands the world
model in a distinct basin shaped by the source task---useful context for
interpreting transfer beyond the scalar return, and a reminder that the
reconstruction-free R2 latent is structured enough to read locomotion phase
from directly.

\paragraph{Why R2 / reconstruction-free here.} Replacing the pixel decoder with
a redundancy-reduction objective~\cite{barlow,dreamerpro} keeps the model
focused on control-relevant latent structure and removes the cost and artifacts
of pixel reconstruction; the trade-off is the loss of decoded
``imagination'' visualizations. The clean convergence of the Barlow term and the
strong control performance indicate the objective is adequate for this task,
though this report does not isolate its contribution against the standard
\texttt{dreamer} reconstruction loss.

\section{Threats to validity and caveats}
\label{sec:caveats}

\begin{itemize}
\item \textbf{Single seed.} All three runs are seed~0; there are no error bars.
      The mid-training $\sim\!4\times$ gap and the $+179$ final advantage are
      large relative to typical DMC run-to-run noise, but statistical claims
      would require additional seeds.
\item \textbf{Fixed budget, scratch not converged.} The scratch baseline is
      still rising at 400k steps, so the advantage measures sample efficiency at
      a fixed budget, not an asymptotic ceiling difference.
\item \textbf{Coarse evaluation.} Evaluation uses 5 deterministic episodes every
      50k steps; the curves are therefore mildly noisy and the reported scores
      are point estimates.
\item \textbf{PCA variance.} The latent projections explain only $\sim\!17\%$ of
      the variance in three components; the separation and ring structure are
      real on the dominant axes but should not be read as the full latent
      geometry.
\item \textbf{No representation-loss ablation.} We run only
      \texttt{rep\_loss=r2dreamer}; we do not compare against the
      \texttt{dreamer}, \texttt{infonce}, or \texttt{dreamerpro} objectives on
      this task, so the report characterizes \emph{this} configuration rather
      than attributing the result to the R2 objective specifically.
\item \textbf{Decoder-free.} With \texttt{rep\_loss=r2dreamer} there is no
      decoder, so no reconstruction or imagination-quality metric is
      available; visual evidence is limited to executed-policy rollouts and latent
      geometry.
\item \textbf{Interrupted source run.} The source run was resumed
      (three appended metric segments); the analysis uses the final continuous
      segment. The reloaded checkpoint and all transfer results are unaffected.
\end{itemize}

\section{Conclusion and future work}

Within a single A100 budget, R2-Dreamer exhibits clean, sizable positive
transfer on a pixel-based \texttt{walker\_walk}\,$\rightarrow$\,\texttt{walker\_run}
shift---large warm start, several-fold mid-training sample efficiency, a $+179$
final advantage at equal budget---together with a distinct, better-structured
latent representation in the fine-tuned model. The most valuable next steps are:
(1) \textbf{2--3 additional seeds} for the target pair to put error bars on the
advantage; (2) \textbf{extended budgets} for scratch (and fine-tune) to separate
``head start'' from ``higher ceiling''; (3) a \textbf{representation-loss
ablation} (\texttt{dreamer} vs.\ \texttt{r2dreamer} vs.\ \texttt{infonce} vs.\
\texttt{dreamerpro}) to attribute the effect; and (4) \textbf{harder shifts}
(e.g.\ Meta-World task transfer, or proprio$\rightarrow$vision) to test the
generality of the transfer.

\section*{Reproduction}

Open \texttt{notebooks/01\_r2dreamer\_foundation.ipynb} on a Colab A100 with a
Python~3.11 runtime; the default \texttt{a100\_vision} preset reproduces the runs
above (set \texttt{R2\_PRESET=t4\_proprio} for the smaller proprioceptive
variant). After the smoke test, run the source, fine-tune, and scratch cells.
Results, rollout videos, and the latent analysis are regenerated by
\texttt{notebooks/06\_r2dreamer\_results.ipynb}, from which every figure and
number in this report is taken.

\begin{thebibliography}{99}

\bibitem{dreamerv3} D.~Hafner, J.~Pasukonis, J.~Ba, and T.~Lillicrap.
\emph{Mastering Diverse Domains through World Models}. arXiv:2301.04104, 2023.

\bibitem{dreamerv2} D.~Hafner, T.~Lillicrap, M.~Norouzi, and J.~Ba.
\emph{Mastering Atari with Discrete World Models}. ICLR, 2021.
arXiv:2010.02193.

\bibitem{dreamer} D.~Hafner, T.~Lillicrap, J.~Ba, and M.~Norouzi.
\emph{Dream to Control: Learning Behaviors by Latent Imagination}. ICLR, 2020.
arXiv:1912.01603.

\bibitem{planet} D.~Hafner, T.~Lillicrap, I.~Fischer, R.~Villegas, D.~Ha,
H.~Lee, and J.~Davidson. \emph{Learning Latent Dynamics for Planning from
Pixels}. ICML, 2019. arXiv:1811.04551.

\bibitem{dreamerpro} F.~Deng, I.~Jang, and S.~Ahn. \emph{DreamerPro:
Reconstruction-Free Model-Based Reinforcement Learning with Prototypical
Representations}. ICML, 2022. arXiv:2110.14565.

\bibitem{barlow} J.~Zbontar, L.~Jing, I.~Misra, Y.~LeCun, and S.~Deny.
\emph{Barlow Twins: Self-Supervised Learning via Redundancy Reduction}. ICML,
2021. arXiv:2103.03230.

\bibitem{cpc} A.~van den Oord, Y.~Li, and O.~Vinyals.
\emph{Representation Learning with Contrastive Predictive Coding}.
arXiv:1807.03748, 2018.

\bibitem{dmc} Y.~Tassa, Y.~Doron, A.~Muldal, T.~Erez, Y.~Li, D.~de~Las~Casas,
D.~Budden, A.~Abdolmaleki, J.~Merel, A.~Lefrancq, T.~Lillicrap, and
M.~Riedmiller. \emph{DeepMind Control Suite}. arXiv:1801.00690, 2018.

\bibitem{dmcontrol} S.~Tunyasuvunakool, A.~Muldal, Y.~Doron, S.~Liu, S.~Bohez,
J.~Merel, T.~Erez, T.~Lillicrap, N.~Heess, and Y.~Tassa.
\emph{dm\_control: Software and Tasks for Continuous Control}.
Software Impacts, 6:100022, 2020.

\bibitem{mujoco} E.~Todorov, T.~Erez, and Y.~Tassa.
\emph{MuJoCo: A Physics Engine for Model-Based Control}. IROS, 2012.

\bibitem{r2dreamer} NM512. \emph{r2dreamer}. GitHub repository,
\url{https://github.com/NM512/r2dreamer}.

\bibitem{dreamerv3torch} NM512. \emph{dreamerv3-torch: PyTorch Implementation of
DreamerV3}. GitHub repository, \url{https://github.com/NM512/dreamerv3-torch}.

\end{thebibliography}

\end{document}
